\subsection{同伦}

定义4.——设(C, d)和(C', d')为两个复形，并设f和g为从C到C'的两个态射。连接 f 到 g 的一个同伦是指从C到C'的任意次数为1的分次A-同态s，使得 \(g - f = d' \circ s + s \circ d\) .

我们称 f 与 g 同伦，如果存在一个连接 f 到 g 的同伦。

若 \(h\) 是从 \(C\) 到 \(C'\) 的第三个态射，且 \(s\) （相应地 \(t\) ）是连接 \(f\) 到 \(g\) （相应地 \(g\) 到 \(h\) ）的一个同伦，则 \(s + t\) 是连接 \(f\) 到 \(h\) 的一个同伦；因此，关系“f 与 g 是从 C 到 C' 的两个同伦的态射”是一个等价关系，其等价类称为从 C 到 C' 的态射的同伦类。

给定两个拓扑空间 X 和 Y，以及一个连续映射 \(f: \mathbf{X} \to \mathbf{Y}\) ，我们定义一个线性映射 \(f_{*}\) 从奇异复形（参见第 3 号）C(X, A) 到 C(Y, A)，即设定 \(f_{*}(e_{s}) = e_{f \circ s}\) （其中 \(s \in \Sigma_{n}(\mathbf{X})\) ）。这个映射是一个复形态射。

两个连续映射 \(f\) 与 \(g\) （从 X 到 Y）称为拓扑同伦的，如果存在一个连续映射 \(h\) ，它把 \([0, 1] \times X\) 映到 Y 中，使得 \(h(0, x) = f(x)\) 与 \(h(1, x) = g(x)\) （对所有 \(x \in X\) ）。我们证明，如果 \(f\) 与 \(g\) 是拓扑同伦的，则态射 \(f_*\) 与 \(g_*\) 在上述定义 4 的意义上是同伦的。正是这一事实构成了代数中所用术语的来源。*

命题3. --- 若f和g是从C到C'的两个同伦的态射，则H(f) = H(g)。

设 s 是从 f 到 g 的同伦。我们有

\[
(g - f) (\mathbf {Z} (\mathbf {C})) = (d ^ {\prime} \circ s + s \circ d) (\mathbf {Z} (\mathbf {C})) = (d ^ {\prime} \circ s) (\mathbf {Z} (\mathbf {C})) \subset \mathrm{B} (\mathbf {C} ^ {\prime}),
\]

因此 \(\mathrm{H}(g - f) = 0\) 且 \(\mathrm{H}(g) = \mathrm{H}(f)\) .

推论。--- 与拟同构同伦的态射是拟同构。

命题4. --- 设 C, C', D, D' 为四个复形， \(f: \mathbf{C} \to \mathbf{C}'\) , \(g: \mathbf{C} \to \mathbf{C}'\) , \(u: \mathbf{D} \to \mathbf{C}\) , \(v: \mathbf{C}' \to \mathbf{D}'\) 四个态射。若 s 是连接 \(f\) 到 \(g\) 的一个同伦，则 \(v \circ s \circ u\) 是连接 \(v \circ f \circ u\) 到 \(v \circ g \circ u\) 的一个同伦。如果 f 和 g 是同伦的，那么 \(v \circ f \circ u\) 与 \(v \circ g \circ u\) 也是同伦的 .

这是显然的。

推论。——设C、C′、C″为三个复形，f和g为从C到C′的两个态射， \(f_{1}\) 与 \(g_{1}\) 为从 C' 到 C'' 的两个态射。若 s 和 \(s_{1}\) 分别是连接 f 到 g 与连接 \(f_{1}\) 到 \(g_{1}\) 的同伦，则 \(s_{1} \circ f + g_{1} \circ s\) 是连接 \(f_{1} \circ f\) 到 \(g_{1} \circ g\) 的一个同伦。如果 f 与 \(f_{1}\) 分别同伦于 g 与 \(g_{1}\) ，则 \(f_{1} \circ f\) 同伦于 \(g_{1} \circ g\) .

事实上， \(s_{1} \circ f\) 连接 \(f_{1} \circ f\) 到 \(g_{1} \circ f\) ，而 \(g_{1} \circ s\) 连接 \(g_{1} \circ f\) 到 \(g_{1} \circ g\) .

定义 5. --- 复形的态射 \(f: \mathbb{C} \to \mathbb{C}'\) 称为同伦等价，如果存在一个态射 \(f': \mathbb{C}' \to \mathbb{C}\) 使得 \(f' \circ f\) 与 \(f \circ f'\) 分别同伦于 \(1_{\mathbb{C}}\) 与 \(1_{\mathbb{C}'}\) 。

显然， \(f'\) 因此也是一个同伦等价；人们也说 \(f'\) 是 \(f\) 在同伦意义下的逆。如果 \(f'\) 与 \(f_1'\) 都是 \(f\) 在同伦意义下的逆，则 \(f'\) 与 \(f_1'\) 是同伦的（事实上，由前述推论， \(f_1' = f_1' \circ 1_{\mathbf{C}'}\) 同伦于 \(f_1' \circ f \circ f'\) ，从而同伦于 \(1_{\mathbf{C}} \circ f' = f'\) ).

命题5. — 同伦等价是拟同构；由同伦等价复合而成的态射是同伦等价。与同伦等价同伦的态射是同伦等价。

设 \(f: \mathbf{C} \to \mathbf{C}'\) 与 \(f_1: \mathbf{C}' \to \mathbf{C}''\) 是复形的同伦等价， \(f': \mathbf{C}' \to \mathbf{C}\) 与 \(f_1': \mathbf{C}'' \to \mathbf{C}'\) 是同伦意义下互逆的态射。我们有

\[
\mathrm{H} (f ^ {\prime}) \circ \mathrm{H} (f) = \mathrm{H} (f ^ {\prime} \circ f) = \mathrm{H} (1 _ {\mathrm{C}}) = 1 _ {\mathrm{H(C)}} \quad (\text { prop. } 3)
\]

并且类似地 \(\mathrm{H}(f) \circ \mathrm{H}(f') = 1_{\mathrm{H}(\mathbb{C}')}\) ，因此 \(\mathrm{H}(f)\) 是双射且 \(f\) 是一个拟同构。另一方面， \((f' \circ f_1') \circ (f_1 \circ f)\) 同伦于 \(f' \circ 1_{\mathbb{C}'} \circ f\) （命题4），从而同伦于 \(1_{\mathbb{C}}\) ；同样地， \((f_1 \circ f) \circ (f' \circ f_1')\) 同伦于 \(1_{\mathbb{C}''}\) ，且 \(f_1 \circ f\) 是一个同伦等价。最后，如果 \(g: \mathbb{C} \to \mathbb{C}'\) 是同伦于 \(f\) 的一个态射，则 \(f' \circ g\) 同伦于 \(f' \circ f\) ，从而同伦于 \(1_{\mathbb{C}}\) ；同样地， \(g \circ f'\) 同伦于 \(f \circ f'\) ，从而同伦于 \(1_{\mathbb{C}'}\) ，且 \(g\) 是一个同伦等价。

推论。——设C、C'、D、D'为四个复形， \(f: \mathbf{C} \to \mathbf{C}'\) 一个态射， \(u: \mathbf{D} \to \mathbf{C}\) 与 \(v: \mathbf{C}' \to \mathbf{D}'\) 为同伦等价。要使 \(v \circ f \circ u\) 是同伦等价（相应地，拟同构），必要且充分的条件是 \(f\) 是一个。

若 \(f\) 是一个同伦等价（相应地，拟同构），则 \(v \circ f \circ u\) 由同伦等价（相应地，拟同构）复合而成，因此也是一个。反过来，设 \(\overline{u}\) 与 \(\overline{v}\) 是在同伦意义下逆于 \(u\) 与 \(v\) 的态射；则 \(\overline{v} \circ (v \circ f \circ u) \circ \overline{u}\) 由命题4同伦于 \(f\) ；因此由命题5及命题3的推论得出结论。

我们称复形C同伦于零，如果 \(1_{C}\) 同伦于零映射，即，若存在 C 的一个次数为 1 的分次 A-自同态 s，使得 \(1_{C} = s \circ d + d \circ s\) 。这也等价于说唯一态射 \(0 \to C\) (相应地， \(C \to 0\) )是一个同伦等价。一个同伦于零的复形具有零同调（命题5）。

例. --- 设 \(u: \mathbf{M} \to \mathbf{N}\) 与 \(v: \mathbf{N} \to \mathbf{P}\) 是 A-模的同态，使得 \(v \circ u = 0\) ；令 C 为满足下述条件的复形： \(\mathbf{C}_2 = \mathbf{M}\) , \(\mathbf{C}_1 = \mathbf{N}\) , \(\mathbf{C}_0 = \mathbf{P}\) , \(\mathbf{C}_i = 0\) 当 \(i \neq 0, 1, 2, d_2 = u, d_1 = v, d_i = 0\) 当 \(i \neq 1, 2\) 。则 C 的同调为零当且仅当序列 \(0 \to \mathbf{M} \xrightarrow{u} \mathbf{N} \xrightarrow{v} \mathbf{P} \to 0\) 是正合的。它同伦于零当且仅当该序列是分裂的。事实上，说 C 同伦于零意味着存在 A-同态 \(s: \mathbf{P} \to \mathbf{N}\) 与 \(t: \mathbf{N} \to \mathbf{M}\) 使得 \(v \circ s = 1_{\mathbf{P}}, s \circ v + u \circ t = 1_{\mathbf{N}}, t \circ u = 1_{\mathbf{M}}\) ；这蕴含该序列是分裂的；反之，若 \(s\) 是 \(v\) 的一个 A-线性截面，我们如下定义 \(t\) ： \(u \circ t = 1_{\mathbf{N}} - s \circ v\) ，这之所以可能，是因为 \(v \circ (1_{\mathbf{N}} - s \circ v) = v - v \circ s \circ v = 0\) .

